\subsection{沿闭子集的深度}

设 A 为 Noether 环，F 为 \(\operatorname{Spec}(A)\) 的闭子集，M 为 A-模。由第 1 节的命题 2 的推论 2，元素 \(\operatorname{prof}_{A}(J;M)\) （属于 \(\mathbf{N} \cup \{+\infty\}\) ）不依赖于满足 \(F = V(J)\) 的 A 的理想 J；它称为 M 沿 F 的深度，记作 \(\operatorname{prof}_{F}(M)\) 。

注.-- 1) 设 \(r\) 为整数。由第 1 节的命题 2 及 II, § 4, 第 4 节, 命题 17 的推论 2，不等式 \(\mathrm{prof}_{\mathrm{F}}(\mathbf{M}) \geqslant r\) 等价于下述性质：对每个支撑集包含于 F 的有限生成 A-模 N，\(\mathrm{Ext}_{\mathrm{A}}^{i}(\mathrm{N},\mathrm{M}) = 0\) 对 \(i < r\) 成立。

\begin{enumerate}
\def\labelenumi{\arabic{enumi})}
\setcounter{enumi}{1}
\tightlist
\item
  设 A-模 M 有限生成。由第 1 节的注 2 与第 4 节的定理 2，我们有下列等价关系
\end{enumerate}

\[
\begin{array}{c} \operatorname{prof} _ {\mathrm{F}} (\mathrm{M}) = 0 \Longleftrightarrow \operatorname{Ass} (\mathrm{M}) \cap \mathrm{F} \neq \emptyset \\ \operatorname{prof} _ {\mathrm{F}} (\mathrm{M}) <   + \infty \Longleftrightarrow \operatorname{Supp} (\mathrm{M}) \cap \mathrm{F} \neq \emptyset . \end{array}
\]

命题 8. 设 A 为 Noether 环，F 为 \(\operatorname{Spec}(\Lambda)\) 的闭子集，M 为有限生成 A-模。我们有

\[
\operatorname{prof} _ {F} (M) = \inf _ {\mathfrak {p} \in F} \operatorname{prof} _ {A _ {\mathfrak {p}}} \left(M _ {\mathfrak {p}}\right) = \inf _ {\mathfrak {p} \in \operatorname{Supp} (M) \cap F} \operatorname{prof} _ {A _ {\mathfrak {p}}} \left(M _ {\mathfrak {p}}\right).
\]

若 \(\operatorname{prof}_{F}(M) = +\infty\) 这是显然的（注 2）。若 \(\operatorname{prof}_{F}(M) = 0\) ，则存在素理想 \(\mathfrak{p} \in \operatorname{Ass}(M) \cap F\) （注 2）；我们有 \(\mathfrak{p}A_{\mathfrak{p}} \in \operatorname{Ass}(M_{\mathfrak{p}})\) (IV, § 1, 第 2 节, 命题 5)，从而 \(\operatorname{prof}_{\Lambda_{\mathfrak{p}}}(\mathrm{M}_{\mathfrak{p}}) = 0\) （第 1 节的注 2），命题在此情形得证。

设 \(0 < \operatorname{prof}_{\mathrm{F}}(\mathrm{M}) < +\infty\) 。设 J 为 A 的理想使 \(\mathrm{V}(\mathrm{J}) = \mathrm{F}\) ，并设 \(x\) 为 J 中使乘以 x 的映射 \(x_{\mathrm{M}}\) 为单射的元素（同上引文）。对每个素理想 \(\mathfrak{p}\) ，乘以 x 的映射 \(x_{\mathrm{M}_{\mathfrak{p}}}\) 是单射的。于是由第 4 节的命题 7，我们有

\[
\begin{array}{c} \operatorname{prof} _ {\mathrm{F}} (\mathrm{M} / x \mathrm{M}) = \operatorname{prof} _ {\mathrm{F}} (\mathrm{M}) - 1 \\ \operatorname{prof} _ {\Lambda_ {\mathfrak {p}}} \big ((\mathrm{M} / x \mathrm{M}) _ {\mathfrak {p}} \big) = \operatorname{prof} _ {\Lambda_ {\mathfrak {p}}} (\mathrm{M} _ {\mathfrak {p}}) - 1. \end{array}
\]

然后对整数 \(\mathrm{prof}_{\mathbb{F}}(\mathbb{M})\) 作归纳即得结论。

注 3. 若 q 是 Supp(M) 的点，则我们有 \(\operatorname{prof}_{\mathrm{A}}(\mathfrak{q};\mathrm{M}) = \inf_{\mathfrak{p}\supseteq q} \operatorname{prof}_{\Lambda_{\mathfrak{p}}}(\mathrm{M}_{\mathfrak{p}})\) 。特别地，有不等式 \(\operatorname{prof}_{\mathrm{A}}(\mathfrak{q};\mathrm{M}) \leqslant \operatorname{prof}_{\mathrm{A}_{\mathfrak{q}}}(\mathrm{M}_{\mathfrak{q}})\) ；当 q 极大时取等号。在一般情形，可能有 \(\operatorname{prof}_{\mathrm{A}}(\mathfrak{q};\mathrm{M}) < \operatorname{prof}_{\Lambda_{\mathfrak{q}}}(\mathrm{M}_{\mathfrak{q}})\) ；也可能有 \(\operatorname{prof}_{\mathrm{A}}(\mathfrak{q};\mathrm{M}) < \inf \operatorname{prof}_{\mathrm{A}_{\mathfrak{m}}}(\mathrm{M}_{\mathfrak{m}})\) ，其中 m 遍历 A 的包含 q 的极大理想组成的集合。例如，设 p 为 A 的一个非极大素理想，它包含 q 且异于 q；令 M = A/p。则有 \(\operatorname{prof}_{\mathrm{A}}(\mathfrak{q};\mathrm{M}) = 0\) ，\(\operatorname{prof}_{\mathrm{A}_{\mathfrak{q}}}(\mathrm{M}_{\mathfrak{q}}) = +\infty\) ，且对 A 的每个极大理想 m 有 \(\operatorname{prof}_{\mathrm{A}_{\mathfrak{m}}}(\mathrm{M}_{\mathfrak{m}}) > 0\) 。

命题 9.--- 设 A 为 Noether 环，M 与 N 为两个有限生成 A-模，F 为 N 的支撑集。则 \(\operatorname{prof}_{\mathbf{F}}(\mathbf{M})\) 是 \(\mathbf{N} \cup \{+\infty\}\) 中由满足下述条件的整数 \(n\) 组成的集合的下确界：\(\operatorname{Ext}_{\Lambda}^{n}(\mathbf{N}, \mathbf{M}) \neq 0\) 。

由注 1，\(\mathrm{Ext}_{\mathrm{A}}^{i}(\mathrm{N},\mathrm{M}) = 0\) 对每个 \(i < \mathrm{prof}_{\mathrm{F}}(\mathrm{M})\) 成立。剩下要证明：若 \(\mathrm{prof}_{\mathrm{F}}(\mathrm{M}) = n < +\infty\) ，则 \(\mathrm{Ext}_{\mathrm{A}}^{n}(\mathrm{N},\mathrm{M}) \neq 0\) 。设 J 为 N 的零化子；有 \(\mathrm{F} = \mathrm{V}(\mathrm{J})\) ，从而 \(\mathrm{prof}_{\mathrm{F}}(\mathrm{M}) = \mathrm{prof}_{\mathrm{A}}(\mathrm{J};\mathrm{M})\) 。由定理 2 的推论 1（ \(n^{\circ}4\) ），存在 M-正则序列 ( \(x_{1},\ldots,x_{n}\) )，其长度为 \(n\) ，由 J 中元素组成，且 A-模 \(\overline{\mathrm{M}} = \mathrm{M}/(x_{1}\mathrm{M}+\ldots+x_{n}\mathrm{M})\) 关于 J 的深度为零。由 A, X, 第 166 页, 命题 9，只需证明 \(\mathrm{Hom}_{\mathrm{A}}(\mathrm{N},\overline{\mathrm{M}})\) 非零。而由命题 8，存在 \(\mathfrak{p} \in \mathrm{Supp}(\mathrm{M}) \cap \mathrm{Supp}(\mathrm{N})\) 使 \(\mathrm{prof}_{\mathrm{A}_{\mathfrak{p}}}(\overline{\mathrm{M}}_{\mathfrak{p}}) = 0\) ，即 \(\mathrm{Hom}_{\mathrm{A}_{\mathfrak{p}}}(\kappa(\mathfrak{p}),\overline{\mathrm{M}}_{\mathfrak{p}}) \neq 0\) 。由于 \(\mathbf{N}_{\mathfrak{p}}\) 非零，\(\kappa(\mathfrak{p})\) -向量空间 \(\mathbf{N}_{\mathfrak{p}} \otimes_{\mathbf{A}_{\mathfrak{p}}} \kappa(\mathfrak{p})\) 非零（Nakayama 引理），其对偶亦然；因此存在一个满 \(\mathbf{A}_{\mathfrak{p}}\) -线性映射，它从 \(\mathbf{N}_{\mathfrak{p}}\) 映到 \(\kappa(\mathfrak{p})\) 。由此得出 \(\mathrm{Hom}_{\mathrm{A}_{\mathfrak{p}}}(\mathbf{N}_{\mathfrak{p}},\overline{\mathbf{M}}_{\mathfrak{p}}) \neq 0\) ，从而 \(\mathrm{Hom}_{\mathrm{A}}(\mathrm{N},\overline{\mathrm{M}}) \neq 0\) (II, § 2, \(n^{\circ}7\) , 命题 19)，此即所要证明的。

注 4.--- 设 A 为 Noether 环，N 为有限生成 A-模。在 \(\mathbf{N}\cup\{+\infty\}\) 中取由整数 \(n\) （使 \(\operatorname{Ext}_{\mathrm{A}}^{n}(\mathrm{N},\mathrm{A})\) 非零）组成的集合的最大下界，人们有时把它称为 N 的级并记作 grade (N)。由命题 9，它也是 A 沿 N 的支撑集的深度，或者（ \(n^{\circ}\) 4, 定理 2）是 N 的零化子的元素组成的 A-正则序列的长度的集合的最小上界。由于对 A 的每个素理想 p，\(N_{p}\) 的零化子等于 \(\operatorname{Ann}(N)_{p}\) (II, § 2, \(n^{\circ}\) 4, 公式 (9))，由 \(n^{\circ}\) 3 的命题 5 可推出等式

\[
\operatorname{grade} (\mathrm{N}) = \inf _ {\mathfrak {p} \in \operatorname{Spec} (\Lambda)} \operatorname{grade} \left(\mathrm{N} _ {\mathfrak {p}}\right) = \inf _ {\mathfrak {m} \in \Omega} \operatorname{grade} \left(\mathrm{N} _ {\mathfrak {m}}\right),
\]

其中 Ω 表示 A 的极大理想组成的集合。

